\chapter{Spaces of Constant Curvature}
\label{chap:dc-spaces-of-constant-curvature}

Faithful transcription of do Carmo, \emph{Riemannian Geometry}. Each numbered
statement keeps do Carmo's number, kind and (name); proofs keep his explicit
cross-references ("by \cref{thm:dc-ch8-2-1}", "from \cref{cor:dc-ch5-2-5}", "Lemma
3.3 of Chap. 7") so the dependency graph is recoverable from the text.

\section{Introduction}

Among the Riemannian manifolds, those with constant sectional curvature are the
most simple. They are related to classical non-Euclidean geometry, which is
historically the first example of a geometric structure different from the
Euclidean. An important property of the spaces with constant curvature is that
they have a sufficiently large number of local isometries (cf. \cref{cor:dc-ch8-2-2});
this property of "free mobility of small triangles" was considered as a postulate
to be satisfied by the non-Euclidean geometries.

When we multiply a Riemannian metric by a positive constant $c$, then its
sectional curvature is multiplied by $1/c$. Therefore, up to a similarity, we can
suppose that the value of the constant sectional curvature of a Riemannian
manifold is $1$, $0$, or $-1$; this hypothesis is assumed without further comment.
So far we have encountered two examples of Riemannian manifolds with constant
sectional curvature $K$: the Euclidean space $\mathbb{R}^n$ with $K\equiv 0$ and
the unit sphere $S^n\subset\mathbb{R}^{n+1}$ with $K\equiv 1$. In this chapter we
introduce a Riemannian manifold, the \emph{hyperbolic space} $H^n$ of dimension $n$,
which has sectional curvature $K\equiv -1$. The manifolds $\mathbb{R}^n$, $S^n$
and $H^n$ are complete and simply connected (cf. Section 3). The main theorem of
this chapter is that these are essentially the only complete, simply connected
Riemannian manifolds with constant sectional curvature. This allows us to reduce
the problem of finding all the complete manifolds of constant curvature to the
problem of determining certain subgroups of the group of isometries of
$\mathbb{R}^n$, $S^n$ and $H^n$ (cf. Section 4).

To prove this, we introduce in Section 2 a slightly more general theorem on the
determination of the metric by means of the curvature. We frequently use the
classical expression \emph{spaces of constant curvature} to designate the Riemannian
manifolds of constant sectional curvature. In the last section we describe the
isometries of the hyperbolic space $H^n$ and identify certain important
hypersurfaces of $H^n$, namely, the horospheres and the hyperspheres.

\section{Theorem of Cartan on the determination of the metric by means of the curvature}

Let $M$ and $\tilde{M}$ be two Riemannian manifolds of dimension $n$ and let
$p\in M$ and $\tilde{p}\in\tilde{M}$. Choose a linear isometry
$i:T_p(M)\to T_{\tilde{p}}(\tilde{M})$. Let $V\subset M$ be a normal neighborhood
of $p$ such that $\exp_{\tilde{p}}$ is defined at $i\circ\exp_p^{-1}(V)$. Define a
mapping $f:V\to\tilde{M}$ by

$$
f(q)=\exp_{\tilde{p}}\circ\,i\circ\exp_p^{-1}(q),\qquad q\in V.
$$

For all $q\in V$ there exists a unique normalized geodesic $\gamma:[0,t]\to M$
with $\gamma(0)=p$, $\gamma(t)=q$. Denote by $P_t$ the parallel transport along
$\gamma$ from $\gamma(0)$ to $\gamma(t)$. Define $\phi_t:T_q(M)\to T_{f(q)}(\tilde{M})$
by

$$
\phi_t(v)=\tilde{P}_t\circ i\circ P_t^{-1}(v),\qquad v\in T_q(M),
$$

where $\tilde{P}_t$ is the parallel transport along the normalized geodesic
$\tilde{\gamma}:[0,t]\to\tilde{M}$ given by $\tilde{\gamma}(0)=\tilde{p}$,
$\tilde{\gamma}'(0)=i(\gamma'(0))$. Denote by $R$ and $\tilde{R}$ the curvatures of
$M$ and $\tilde{M}$, respectively.

\begin{theorem}[E. Cartan]
  \label{thm:dc-ch8-2-1}
  \dcref{ch8:2.1}
  \uses{cor:dc-ch5-2-5}
  \notready
  With the notation above, if for all $q\in V$ and all $x,y,u,v\in T_q(M)$ we have
  
  $$
  \langle R(x,y)u,v\rangle=\big\langle\tilde{R}(\phi_t(x),\phi_t(y))\phi_t(u),\phi_t(v)\big\rangle,
  $$
  
  then $f:V\to f(V)\subset\tilde{M}$ is a local isometry and $df_p=i$.
\end{theorem}
\begin{proof}
  Let $q\in V$ and let $\gamma:[0,\ell]\to M$ be a normalized geodesic with
  $\gamma(0)=p$, $\gamma(\ell)=q$. Let $v\in T_q(M)$ and let $J$ be a Jacobi field
  along $\gamma$ with $J(0)=0$ and $J(\ell)=v$. Let $e_1,\dots,e_n=\gamma'(0)$ be an
  orthonormal basis of $T_p(M)$ and let $e_i(t)$ be the parallel transport of $e_i$
  along $\gamma$. Writing $J(t)=\sum_i y_i(t)e_i(t)$, the Jacobi equation gives
  
  $$
  y_j''+\sum_i\langle R(e_n,e_i)e_n,e_j\rangle y_i=0,\qquad j=1,\dots,n.
  $$
  
  Let $\tilde{\gamma}:[0,\ell]\to\tilde{M}$ be the normalized geodesic with
  $\tilde{\gamma}(0)=\tilde{p}$, $\tilde{\gamma}'(0)=i(\gamma'(0))$, and set
  $\tilde{J}(t)=\phi_t(J(t))$, $\tilde{e}_j(t)=\phi_t(e_j(t))$. By linearity of
  $\phi_t$, $\tilde{J}(t)=\sum_i y_i(t)\tilde{e}_i(t)$. Since by hypothesis
  $\langle R(e_n,e_i)e_n,e_j\rangle=\langle\tilde{R}(\tilde{e}_n,\tilde{e}_i)\tilde{e}_n,\tilde{e}_j\rangle$,
  $\tilde{J}$ satisfies the same equation, hence $\tilde{J}$ is a Jacobi field along
  $\tilde{\gamma}$ with $\tilde{J}(0)=0$. Since parallel transport is an isometry,
  $|\tilde{J}(\ell)|=|J(\ell)|$. It remains to show that $\tilde{J}(\ell)=df_q(v)=df_q(J(\ell))$.
  Since $\tilde{J}(t)=\phi_t(J(t))$, $\tilde{J}'(0)=i(J'(0))$. Since $J$ and
  $\tilde{J}$ are Jacobi fields vanishing at $t=0$, by \cref{cor:dc-ch5-2-5},
  
  $$
  J(t)=(d\exp_p)_{t\gamma'(0)}(tJ'(0)),\qquad\tilde{J}(t)=(d\exp_{\tilde{p}})_{t\tilde{\gamma}'(0)}(t\tilde{J}'(0)).
  $$
  
  Therefore
  
  $$
  \tilde{J}(\ell)=(d\exp_{\tilde{p}})_{\ell\tilde{\gamma}'(0)}\circ i\circ\big((d\exp_p)_{\ell\gamma'(0)}\big)^{-1}(J(\ell))=df_q(J(\ell)),
  $$
  
  which proves the assertion. $\square$
  
  Observe that the same proof shows that if $\exp_p$ and $\exp_{\tilde{p}}$ are
  diffeomorphisms, then, under the conditions of Theorem 2.1, $f$ is defined on all
  of $M$ and is an isometry. The theorem implies that the metric is, in a certain
  sense, determined locally by the curvature. An equivalent assertion was made by
  Riemann; the first proof of the local theorem was presented by E. Cartan. A global
  version was given by W. Ambrose in 1956.
\end{proof}

\begin{corollary}[Corollary 2.2]
  \label{cor:dc-ch8-2-2}
  \dcref{ch8:2.2}
  \uses{thm:dc-ch8-2-1}
  Let $M$ and $\tilde{M}$ be spaces with the same constant curvature and the same
  dimension $n$. Let $p\in M$ and $\tilde{p}\in\tilde{M}$. Choose arbitrary
  orthonormal bases $\{e_j\}$ of $T_p(M)$ and $\{\tilde{e}_j\}$ of $T_{\tilde{p}}(\tilde{M})$,
  $j=1,\dots,n$. Then there exist a neighborhood $V\subset M$ of $p$, a neighborhood
  $\tilde{V}\subset\tilde{M}$ of $\tilde{p}$, and an isometry $f:V\to\tilde{V}$ such
  that $df_p(e_j)=\tilde{e}_j$.
\end{corollary}
\begin{proof}
  Choose the isometry $i$ of the theorem in such a way that
  $i(e_j)=\tilde{e}_j$. The condition on the curvature is immediately verified, and
  the conclusion follows from \cref{thm:dc-ch8-2-1}. $\square$
\end{proof}

\begin{corollary}[Corollary 2.3]
  \label{cor:dc-ch8-2-3}
  \dcref{ch8:2.3}
  Let $M$ be a space of constant curvature and let $p$ and $q$ be any two points of
  $M$. Let $\{e_j\}$ and $\{f_j\}$ be arbitrary orthonormal bases of $T_p(M)$ and
  $T_q(M)$, respectively. Then there exist neighborhoods $U$ of $p$ and $V$ of $q$,
  and an isometry $g:U\to V$ such that $dg_p(e_j)=f_j$.
  
  Related to the above is the problem of deciding whether a diffeomorphism
  $f:M\to\tilde{M}$ which preserves the curvature in the sense that
  
  $$
  \langle R(X,Y)Z,T\rangle_p=\big\langle\tilde{R}(df_p(X),df_p(Y))df_p(Z),df_p(T)\big\rangle_{f(p)}
  $$
  
  for all $p\in M$ and all $X,Y,Z,T\in T_p(M)$ is an isometry. In dimension two this
  is a kind of converse to the Theorem Egregium of Gauss, and is false, even in the
  compact case (Fig. 1). For $n=\dim M=\dim\tilde{M}\geq 4$, the problem
  surprisingly admits an affirmative solution (Kulkarni); for $n=3$, the problem was
  treated by Yau.
\end{corollary}

\section{Hyperbolic space}

We give now an example of a space of constant curvature $-1$. Consider the
half-space of $\mathbb{R}^n$ given by

$$
H^n=\{(x_1,\dots,x_n)\in\mathbb{R}^n;\ x_n>0\}
$$

and introduce on $H^n$ the metric

$$
g_{ij}(x_1,\dots,x_n)=\frac{\delta_{ij}}{x_n^2}.\qquad(1)
$$

It is clear that $H^n$ is simply connected. We are going to show that $H^n$ with
the metric $(1)$ has constant sectional curvature equal to $-1$, and that $H^n$ is
complete. $H^n$ is called the \emph{hyperbolic space} of dimension $n$.

\begin{definition}[hyperbolic space $H^n$ as a Riemannian manifold]
  \label{def:dc-ch8-3-hyperbolic-space}
  \lean{Riemannian.opensConformalMetric,
    Riemannian.Hyperbolic.hyperbolicMetric,
    Riemannian.Hyperbolic.hyperbolicMetric_apply,
    Riemannian.Hyperbolic.hyperbolicMetric_gram}
  \leanok
  The \emph{hyperbolic space} $H^n$ is the open half-space
  $H^n=\{x\in\mathbb R^n : x_n>0\}$ carrying the metric $g_{ij}=\delta_{ij}/x_n^2$
  of $(1)$, that is, the conformal rescaling of the ambient Euclidean metric by
  the positive smooth factor $1/x_n^2$: for $u,v\in T_pH^n$,
  $\langle u,v\rangle_g=\tfrac{1}{x_n^2}\langle u,v\rangle$. This is a Riemannian
  manifold. More generally, given an open subset of a (complete) real
  inner-product space and a positive smooth function $\varphi$ on it, the
  pointwise rescaling $\langle u,v\rangle\mapsto\varphi\,\langle u,v\rangle$ of
  the ambient Euclidean metric is a Riemannian metric — the conformal metric with
  conformal factor $\varphi$ (positive-definiteness and von-Neumann boundedness
  are inherited from the flat metric, whose metric ball is a Euclidean ball, and
  smoothness of the metric section from that of $\varphi$).
\end{definition}

A good part of the calculation can be carried out in a more general situation. Two
metrics $\langle\,,\rangle$ and $\langle\!\langle\,,\rangle\!\rangle$ on a
differentiable $M$ are \emph{conformal} if there exists a positive differentiable
function $f:M\to\mathbb{R}$ such that
$\langle u,v\rangle_p=f(p)\langle\!\langle u,v\rangle\!\rangle_p$ for all $p\in M$
and $u,v\in T_p(M)$. The metric $(1)$ of $H^n$ is conformal to the usual metric of
Euclidean space. Consider on $H^n$ the metric $g_{ij}=\delta_{ij}/F^2$, where $F$
is a positive differentiable function on $H^n$; write $g^{ij}=F^2\delta_{ij}$ for
the inverse matrix, and put $\log F=f$, $\partial f/\partial x_j=f_j$. Then

$$
\frac{\partial g_{ik}}{\partial x_j}=-2\frac{\delta_{ik}}{F^2}f_j.
$$

A computation of the Christoffel symbols gives

$$
\Gamma_{ij}^k=-\delta_{jk}f_i-\delta_{ki}f_j+\delta_{ij}f_k;
$$

hence if all three indices are distinct, $\Gamma_{ij}^k=0$, while if two indices
are equal,

$$
\Gamma_{ij}^i=-f_j,\quad\Gamma_{ii}^j=f_j,\quad\Gamma_{ij}^j=-f_i,\quad\Gamma_{ii}^i=-f_i.
$$

A direct calculation of the curvature coefficients yields

$$
F^2 R_{ijij}=-\sum_\ell f_\ell^2+f_i^2+f_j^2+f_{ii}+f_{jj},
$$

and $R_{ijk\ell}=0$ if all four indices are distinct; if any three indices are
distinct,

$$
R_{ijk}^i=-f_k f_j-f_{kj},\quad R_{ijk}^j=f_i f_k+f_{ki},\quad R_{ijk}^k=0.\qquad(2)
$$

The sectional curvature with respect to the plane generated by
$\partial/\partial x_i,\partial/\partial x_j$ (which are orthogonal) is

$$
K_{ij}=\frac{R_{ijij}}{g_{ii}g_{jj}}=R_{ijij}F^4=\Big(-\sum_\ell f_\ell^2+f_i^2+f_j^2+f_{ii}+f_{jj}\Big)F^2.
$$

Specializing to $F^2=x_n^2$, so that $f=\log x_n$: if $i\neq n$ and $j\neq n$,
$K_{ij}=-\tfrac{1}{x_n^2}x_n^2=-1$; if $i=n,j\neq n$,
$K_{nj}=(-f_n^2+f_n^2+f_{nn})F^2=-\tfrac{1}{x_n^2}x_n^2=-1$; and similarly if
$i\neq n,j=n$, $K_{in}=-1$. Using $(2)$ and \cref{cor:dc-ch4-3-5}, we conclude
that the sectional curvature of $H^n$ is a constant equal to $-1$.

\begin{lemma}[Christoffel symbols and coordinate curvature of a conformal metric]
  \label{lem:dc-ch8-3-conformal-curvature}
  \lean{Riemannian.Hyperbolic.christoffelFromMetric_eq,
    Riemannian.Hyperbolic.Rcoeff_diag,
    Riemannian.Hyperbolic.hyperbolic_sectionalCurvature}
  \leanok
  For the conformal metric $g_{ij}=\delta_{ij}/F^2$ ($f=\log F$), the Christoffel
  symbols computed from the Koszul formula
  $\Gamma^k_{ij}=\tfrac12\sum_\ell g^{k\ell}(\partial_i g_{j\ell}+\partial_j g_{i\ell}-\partial_\ell g_{ij})$
  equal $\Gamma^k_{ij}=-\delta_{jk}f_i-\delta_{ki}f_j+\delta_{ij}f_k$, and the
  associated coordinate curvature coefficients satisfy, for $i\neq j$,
  $$F^2R_{ijij}=-\sum_\ell f_\ell^2+f_i^2+f_j^2+f_{ii}+f_{jj}.$$
  Specializing to $F=x_n$ (so $f=\log x_n$, $f_a=\delta_{an}/x_n$,
  $f_{ab}=-\delta_{an}\delta_{bn}/x_n^2$) gives, for every \emph{coordinate}
  $2$-plane spanned by $\partial/\partial x_i,\partial/\partial x_j$ with $i\neq j$,
  the coordinate-frame sectional curvature $K_{ij}=R_{ijij}F^4=-1$. (This is the
  coordinate identity; extending it to all $2$-planes and to the intrinsic
  curvature of the manifold is \cref{prop:dc-ch8-3-const-curv}.)
\end{lemma}
\begin{proof}
  \leanok
  The Koszul identity is a direct computation from
  $g^{k\ell}=\delta_{k\ell}F^2$ and $\partial_m g_{ab}=-2\delta_{ab}f_m/F^2$.
  The curvature coefficient is
  $R^s_{ijk}=\partial_j\Gamma^s_{ik}-\partial_i\Gamma^s_{jk}+\sum_\ell(\Gamma^\ell_{ik}\Gamma^s_{j\ell}-\Gamma^\ell_{jk}\Gamma^s_{i\ell})$;
  substituting the closed form for $\Gamma$ and contracting the Kronecker deltas
  yields $R^j_{iji}=-\sum_\ell f_\ell^2+f_i^2+f_j^2+f_{ii}+f_{jj}$, which is
  $F^2R_{ijij}$ since $g_{jj}=1/F^2$. For $F=x_n$ one has
  $(f_a)^2+f_{aa}=0$ at every index, so the bracketed sum collapses to
  $-\sum_\ell f_\ell^2=-1/x_n^2$, and $K_{ij}=R_{ijij}F^4=(F^2R_{ijij})F^2=-1$.
\end{proof}

\begin{lemma}[the Riemann curvature tensor in chart coordinates]
  \label{lem:dc-ch8-3-chart-curvature}
  \lean{Riemannian.leviCivita_curvature_chartFrame_expansion}
  \leanok
  Let $g$ be a Riemannian metric on $M$ and $\nabla$ its Levi-Civita connection.
  In the coordinate frame $\partial_i$ of a chart at a point $p$, the intrinsic
  curvature operator is computed from the chart Christoffel symbols
  $\Gamma^k_{ij}$ by the classical coordinate formula
  $$
  R(\partial_i,\partial_j)\partial_k
    = \sum_\ell\Big(\partial_j\Gamma^\ell_{ik}-\partial_i\Gamma^\ell_{jk}
      +\sum_s(\Gamma^s_{ik}\Gamma^\ell_{js}-\Gamma^s_{jk}\Gamma^\ell_{is})\Big)\partial_\ell,
  $$
  where $\partial_j\Gamma^\ell_{ik}$ is the chart partial derivative of the
  coordinate Christoffel function. With do Carmo's sign convention
  $R(X,Y)Z=\nabla_Y\nabla_X Z-\nabla_X\nabla_Y Z+\nabla_{[X,Y]}Z$ the bracketed
  coefficient is exactly the $R^\ell_{ijk}$ of $(2)$.
\end{lemma}
\begin{proof}
  \leanok
  In a coordinate chart the frame vectors have vanishing pairwise Lie brackets,
  $[\partial_i,\partial_j]=0$, so the bracket term
  $\nabla_{[\partial_i,\partial_j]}\partial_k$ drops out. The field
  $\nabla_{\partial_j}\partial_k$ is, throughout the chart, the frame combination
  $\sum_m\Gamma^m_{jk}\partial_m$ (formula $(10)$ of Chap.~2, the Koszul collapse
  on the coordinate frame, valid at every chart point). Differentiating this
  along $\partial_i$ by the Leibniz rule for $\nabla$ gives
  $\nabla_{\partial_i}\nabla_{\partial_j}\partial_k
   =\sum_\ell\big(\partial_i\Gamma^\ell_{jk}
   +\sum_s\Gamma^s_{jk}\Gamma^\ell_{is}\big)\partial_\ell$, using
  $\nabla_{\partial_i}\partial_m=\sum_\ell\Gamma^\ell_{im}\partial_\ell$ for the
  inner derivative. Subtracting the term with $i$ and $j$ interchanged yields the
  stated coefficients.
\end{proof}

\begin{lemma}[hyperbolic chart Gram and inverse-Gram in the abstract chart frame]
  \label{lem:dc-ch8-3-hyperbolic-chart-gram}
  \lean{Riemannian.Hyperbolic.hyperbolic_chartGramMatrix,
    Riemannian.Hyperbolic.hyperbolic_chartGramMatrix_eq,
    Riemannian.Hyperbolic.hyperbolic_chartInvGramMatrix,
    Riemannian.Hyperbolic.finBasisGramMatrix_det_isUnit}
  \leanok
  \uses{def:dc-ch8-3-hyperbolic-space}
  Because $H^n=\{x_n>0\}$ is an open subset of $\mathbb R^n$, its tangent bundle is
  canonically trivial and the abstract chart frame vector
  $\partial_i=$ \emph{finBasis}$_i$ (the frame underlying \cref{lem:dc-ch8-3-chart-curvature})
  is literally the constant basis vector. Hence the chart Gram matrix of the
  hyperbolic metric is the conformal rescaling
  $G_{ij}(x)=x_n^{-2}\,B_{ij}$ of the constant Gram matrix
  $B_{ij}=\langle\text{finBasis}_i,\text{finBasis}_j\rangle$ of the abstract basis,
  and its inverse is $G^{ij}(x)=x_n^{2}\,B^{ij}$ (with $B$ invertible, being the
  Gram matrix of a basis). Since the abstract basis need not be
  $\langle\,,\rangle$-orthonormal, $B$ is a \emph{general} constant symmetric
  positive-definite matrix, not $\delta$: do Carmo's $\delta_{ij}$ of $(1)$ is
  replaced by $B_{ij}$ in the abstract chart frame, and the whole §3 conformal
  computation goes through with $\delta\rightsquigarrow B$, the extra factors
  collapsing by the inverse-Gram identity $\sum_{ij}B^{ij}c_ic_j=1$ (for
  $c_i=\langle e_n,\text{finBasis}_i\rangle$).
\end{lemma}
\begin{proof}
  \leanok
  The trivialisation of the tangent bundle of an open subset of $\mathbb R^n$ is
  the identity on fibres, so $\partial_i(x)=$ finBasis$_i$ and
  $G_{ij}(x)=\langle\partial_i,\partial_j\rangle_g=x_n^{-2}\langle\text{finBasis}_i,\text{finBasis}_j\rangle=x_n^{-2}B_{ij}$.
  The Gram matrix at the basepoint is positive definite, so $\det B$ is a unit and
  the inverse of the scalar multiple $x_n^{-2}B$ is $x_n^{2}B^{-1}$.
\end{proof}

\begin{lemma}[hyperbolic chart Christoffel symbols and their derivatives, closed form]
  \label{lem:dc-ch8-3-hyperbolic-christoffel}
  \lean{Riemannian.Hyperbolic.hyperbolic_chartChristoffel,
    Riemannian.Hyperbolic.hyperbolic_partialDeriv_chartChristoffel,
    Riemannian.Hyperbolic.hyperbolic_partialDeriv_chartGramOnE}
  \leanok
  \uses{lem:dc-ch8-3-hyperbolic-chart-gram}
  In the abstract chart frame the Christoffel symbols of the hyperbolic metric are,
  in closed form,
  $$\Gamma^k_{ij}(y)=-f_i\delta^k_j-f_j\delta^k_i+d^k B_{ij},\qquad
    f_i=\frac{c_i}{y_n},\ c_i=(\text{finBasis}_i)_n,\ d^k=\frac{\sum_\ell B^{k\ell}c_\ell}{y_n},$$
  the constant abstract-basis Gram $B$ (\cref{lem:dc-ch8-3-hyperbolic-chart-gram})
  playing the role of do Carmo's $\delta$
  ($\Gamma^k_{ij}=-\delta_{jk}f_i-\delta_{ki}f_j+\delta_{ij}f_k$ with $\delta\rightsquigarrow B$).
  Because every term carries the single factor $1/y_n$, one may write
  $\Gamma^k_{ij}(y)=K^k_{ij}\,(y_n)^{-1}$ with $K^k_{ij}$ the ($y$-independent)
  numerator, so the abstract-chart-frame partial derivatives are
  $$\partial_m\Gamma^k_{ij}(y)=K^k_{ij}\cdot\bigl(-c_m\,(y_n)^{-2}\bigr).$$
  This is exactly the $\Gamma$ and $\partial\Gamma$ data consumed by the chart-frame
  curvature expansion \cref{lem:dc-ch8-3-chart-curvature}.
\end{lemma}
\begin{proof}
  \leanok
  The chart Gram is $G_{ij}(y)=(y_n^2)^{-1}B_{ij}$ and its inverse
  $G^{k\ell}(y)=y_n^2 B^{k\ell}$ (\cref{lem:dc-ch8-3-hyperbolic-chart-gram}). The
  distinguished coordinate $y_n=y_e$ has derivative
  $\partial_m y_n=(\text{finBasis}_m)_n=c_m$ along $\text{finBasis}_m$, so
  $\partial_m\bigl((y_n^2)^{-1}\bigr)=-2\,y_n^{-3}c_m$ and hence
  $\partial_m G_{ij}=B_{ij}\cdot(-2\,y_n^{-3}c_m)$. Substituting into the Koszul
  formula
  $\Gamma^k_{ij}=\tfrac12\sum_\ell G^{k\ell}(\partial_iG_{j\ell}+\partial_jG_{i\ell}-\partial_\ell G_{ij})$
  and contracting $\sum_\ell B^{k\ell}B_{\ell j}=\delta^k_j$,
  $\sum_\ell B^{k\ell}B_{\ell i}=\delta^k_i$ (from $B^{-1}B=1$) collapses the factors
  $y_n^2\cdot y_n^{-3}=y_n^{-1}$, giving the closed form. Since only the $1/y_n$ factor
  depends on $y$, $\partial_m\Gamma^k_{ij}=K^k_{ij}\,\partial_m\bigl(y_n^{-1}\bigr)
  =K^k_{ij}\,(-c_m y_n^{-2})$.
\end{proof}

\begin{lemma}[inverse-Gram/coordinate identity $\sum_{ij}B^{ij}c_ic_j=1$]
  \label{lem:dc-ch8-3-inverse-gram-coord}
  \lean{Riemannian.Hyperbolic.finBasis_inv_gram_coord}
  \leanok
  \uses{lem:dc-ch8-3-hyperbolic-chart-gram}
  For the constant abstract-basis Gram matrix
  $B_{ij}=\langle\text{finBasis}_i,\text{finBasis}_j\rangle$ with inverse
  $B^{ij}=(B^{-1})_{ij}$, and $c_i=(\text{finBasis}_i)_n$ the $n$-th coordinate of
  the $i$-th basis vector,
  $$\sum_{i,j} B^{ij}\,c_i\,c_j = 1.$$
  Equivalently $c_i=\langle\text{finBasis}_i, w_n\rangle$ for the standard unit
  vector $w_n=\text{basisFun}_n$, and the $B^{-1}$-contraction reconstructs
  $\langle w_n, w_n\rangle=1$. This is exactly the term
  $|\mathrm d\sigma|^2_{g_0}=1$ that makes the §3 conformal cancellation
  $T=\tfrac12\varphi g_0$ (hence $K\equiv-1$) go through with the \emph{general}
  non-orthonormal abstract Gram $B$ in place of $\delta$.
\end{lemma}
\begin{proof}
  \leanok
  Write $w_n=\sum_i a_i\,\text{finBasis}_i$ in the abstract basis (the
  \texttt{repr} coordinates). Then $c_i=\langle\text{finBasis}_i,w_n\rangle=\sum_j
  B_{ij}a_j$, so $\sum_j B^{ij}c_j=(B^{-1}B a)_i=a_i$ and
  $\sum_{ij}B^{ij}c_ic_j=\sum_i a_i c_i=\langle w_n,w_n\rangle=1$.
\end{proof}

\begin{lemma}[the pointwise $(0,4)$ curvature form and its chart-frame value]
  \label{lem:dc-ch8-3-pointwise-curvature-form}
  \lean{Riemannian.AffineConnection.curvatureFormAt,
    Riemannian.AffineConnection.isAlgCurvatureForm_curvatureFormAt,
    Riemannian.leviCivita_curvatureFormAt_chartFrame}
  \leanok
  \uses{lem:dc-ch8-3-chart-curvature}
  The curvature operator $R(X,Y)Z$ is a tensor: its value at $p$ depends only on
  $X(p),Y(p),Z(p)$ (do Carmo Ch.~4, Prop.~2.2 read pointwise). Lowering by the
  metric gives the pointwise $(0,4)$ form $\langle R(x,y)z,t\rangle_g$ on
  $T_pM$, which for the Levi-Civita connection is an \emph{algebraic curvature
  form} (multilinear, antisymmetric in each pair, first Bianchi — do Carmo
  Prop.~2.5). In the chart frame $\partial_a=\text{finBasis}_a$ its value is the
  coordinate curvature coefficient contracted with the chart Gram,
  $$\langle R(\partial_i,\partial_j)\partial_k,\partial_\ell\rangle_g
    = \sum_m R^m_{ijk}\,G_{m\ell},\qquad
    R^m_{ijk}=\partial_j\Gamma^m_{ik}-\partial_i\Gamma^m_{jk}
      +\sum_s(\Gamma^s_{ik}\Gamma^m_{js}-\Gamma^s_{jk}\Gamma^m_{is}),$$
  the metric-lowered avatar of \cref{lem:dc-ch8-3-chart-curvature}. This bridges
  the \emph{intrinsic} manifold curvature to the coordinate Christoffel data, and
  is the input to the all-planes identification \cref{cor:dc-ch4-3-5}: it makes the
  chart coefficients literally the components of an algebraic curvature form on
  $T_pM$, ready for \texttt{ext\_basis}/\texttt{eq\_kronecker\_iff\_const}.
\end{lemma}
\begin{proof}
  \leanok
  Tensoriality is the pointwise reading of the $\mathcal D(M)$-multilinearity of
  $R$: a field vanishing at $p$ decomposes as $\sum_i f_i W_i+\tau$ with
  $f_i(p)=0$ and $\tau$ vanishing near $p$, and homogeneity of $R$ in that slot
  kills both parts, so $R(\cdot)$ against a field vanishing at $p$ vanishes at
  $p$; hence $R(X,Y)Z|_p$ depends only on the values at $p$. The
  algebraic-curvature-form axioms transport slotwise from the field-level
  symmetries (Prop.~2.5, via the Levi-Civita property). The chart-frame value
  pairs the frame expansion of $R(\partial_i,\partial_j)\partial_k$
  (\cref{lem:dc-ch8-3-chart-curvature}) against $\partial_\ell$, distributing the
  metric over the finite sum.
\end{proof}

\begin{proposition}[$H^n$ has constant sectional curvature $-1$]
  \label{prop:dc-ch8-3-const-curv}
  \lean{Riemannian.Hyperbolic.hyperbolic_isConstantCurvature,
    Riemannian.Hyperbolic.hyperbolicMetric_sectionalCurvature_eq_neg_one,
    Riemannian.Hyperbolic.hyperbolic_curvatureFormAt_eq,
    Riemannian.Hyperbolic.hyperbolic_curvatureFormAt_frame}
  \leanok
  \uses{lem:dc-ch8-3-conformal-curvature,lem:dc-ch8-3-chart-curvature,lem:dc-ch8-3-pointwise-curvature-form,lem:dc-ch8-3-hyperbolic-chart-gram,lem:dc-ch8-3-hyperbolic-christoffel,lem:dc-ch8-3-inverse-gram-coord,cor:dc-ch4-3-5,def:dc-ch8-3-hyperbolic-space}
  The hyperbolic space $H^n$, with the metric $g_{ij}=\delta_{ij}/x_n^2$, has
  constant sectional curvature equal to $-1$: for every point and every
  $2$-plane $\sigma\subset T_pH^n$, $K(p,\sigma)=-1$.
\end{proposition}
\begin{proof}
  \leanok
  Work in the abstract chart frame $\partial_a=\text{finBasis}_a$ of $T_pH^n$. By
  \cref{lem:dc-ch8-3-pointwise-curvature-form} the pointwise $(0,4)$ curvature form
  of the Levi-Civita connection is, in this frame,
  $\langle R(\partial_i,\partial_j)\partial_k,\partial_\ell\rangle_g=\sum_m
  R^m_{ijk}G_{m\ell}$, with $R^m_{ijk}=\partial_j\Gamma^m_{ik}-\partial_i\Gamma^m_{jk}
  +\sum_s(\Gamma^s_{ik}\Gamma^m_{js}-\Gamma^s_{jk}\Gamma^m_{is})$ and $G_{m\ell}$
  the chart Gram matrix. Insert the closed forms
  \cref{lem:dc-ch8-3-hyperbolic-christoffel}, $\Gamma^k_{ij}(y)=y_n^{-1}K^k_{ij}$
  and $\partial_m\Gamma^k_{ij}(y)=-c_m\,y_n^{-2}K^k_{ij}$ with the \emph{constant}
  coefficient $K^k_{ij}=-c_i\delta^k_j-c_j\delta^k_i+D^kB_{ij}$
  ($c_i=(\text{finBasis}_i)_n$, $D^k=\sum_\ell B^{k\ell}c_\ell$), and the chart Gram
  \cref{lem:dc-ch8-3-hyperbolic-chart-gram}, $G_{m\ell}=y_n^{-2}B_{m\ell}$ with $B$
  the constant abstract-basis Gram matrix. Every term carries a factor $y_n^{-4}$,
  which cancels against the right-hand side, and the identity reduces to the purely
  finite contraction $\sum_m R^m_{ijk}B_{m\ell}=-(B_{ik}B_{j\ell}-B_{jk}B_{i\ell})$
  in the constants $c,B,D$. This contraction is do Carmo's conformal cancellation
  $T=\tfrac12\varphi g_0$: expanding the Christoffel products and contracting each
  index sum with the ``lowered'' identity $\sum_m K^m_{ab}B_{m\ell}=-c_aB_{b\ell}
  -c_bB_{a\ell}+c_\ell B_{ab}$ (which uses $\sum_m D^mB_{m\ell}=c_\ell$), every
  monomial cancels except $-(B_{ik}B_{j\ell}-B_{jk}B_{i\ell})$, the last step using
  the single nontrivial input $\sum_{ij}B^{ij}c_ic_j=1$
  (\cref{lem:dc-ch8-3-inverse-gram-coord}). Hence
  $\langle R(\partial_i,\partial_j)\partial_k,\partial_\ell\rangle_g
  =-(G_{ik}G_{j\ell}-G_{jk}G_{i\ell})$ on the chart frame. Both the pointwise
  curvature form and $-1$ times the standard curvature form
  $\langle x,z\rangle\langle y,t\rangle-\langle y,z\rangle\langle x,t\rangle$ are
  algebraic curvature forms (\cref{lem:dc-ch8-3-pointwise-curvature-form}) that
  agree on this basis frame, so by the basis-determination step of
  \cref{cor:dc-ch4-3-5} (\texttt{ext\_basis}) they agree on \emph{all} of $T_pH^n$:
  $\langle R(x,y)z,t\rangle_g=-(\langle x,z\rangle_g\langle y,t\rangle_g
  -\langle y,z\rangle_g\langle x,t\rangle_g)$. This is constant curvature $-1$ at the
  field level; dividing the diagonal $\langle R(x,y)x,y\rangle_g=-|x\wedge y|^2$ by
  $|x\wedge y|^2$ gives $K(p,\sigma)=-1$ for every $2$-plane
  $\sigma=\mathrm{span}\{x,y\}$.
\end{proof}

\begin{lemma}[Christoffel contraction of the hyperbolic metric, closed form]
  \label{lem:dc-ch8-3-christoffel-contraction}
  \lean{Riemannian.Hyperbolic.hyperbolic_chartChristoffelContraction_eq}
  \leanok
  \uses{lem:dc-ch8-3-hyperbolic-christoffel,lem:dc-ch8-3-inverse-gram-coord}
  Contracting the closed-form chart Christoffel symbols
  $\Gamma^k_{ij}=-\delta_{jk}f_i-\delta_{ki}f_j+\delta_{ij}f_k$
  ($f_i=(\partial_i)_n/x_n$) of \cref{lem:dc-ch8-3-hyperbolic-christoffel} against
  two tangent vectors $v,w\in T_pH^n$ collapses, via the inverse-Gram/coordinate
  identity \cref{lem:dc-ch8-3-inverse-gram-coord}, to the basis-independent
  Euclidean-valued expression
  $$
  \Gamma(v,w)(y)=\nabla_v w
    =\frac{1}{y_n}\Big(\langle v,w\rangle\, e_n-v_n\,w-w_n\,v\Big),
  $$
  where $y_n$ is the distinguished coordinate of the base point, $v_n=\langle v,e_n\rangle$,
  $e_n$ is the distinguished unit vector, and $\langle\,,\rangle$ is the ambient
  Euclidean inner product. This is do Carmo's covariant derivative of $H^n$ read in
  the trivial chart of the open half-space.
\end{lemma}
\begin{proof}
  \leanok
  In the abstract chart frame the contraction is
  $\Gamma(v,w)(y)=\sum_k\big(\sum_{ij}\Gamma^k_{ij}(y)\,v^i w^j\big)\partial_k$ with
  $v^i,w^j$ the frame coordinates. Substituting the closed form
  \cref{lem:dc-ch8-3-hyperbolic-christoffel} and summing the two Kronecker terms
  gives $-\big(\sum_i (\partial_i)_n v^i\big)w-\big(\sum_j(\partial_j)_n w^j\big)v
  =-v_n w-w_n v$ (using $\sum_i(\partial_i)_n v^i=v_n$), while the $\delta_{ij}f_k$
  term contracts to $\langle v,w\rangle\sum_k(\sum_\ell B^{k\ell}c_\ell)\partial_k
  =\langle v,w\rangle\,e_n$ by \cref{lem:dc-ch8-3-inverse-gram-coord}; the common
  factor $1/y_n$ is carried throughout.
\end{proof}

\begin{lemma}[geodesic equation of $H^n$ in Euclidean form]
  \label{lem:dc-ch8-3-geodesic-equation}
  \lean{Riemannian.Hyperbolic.hyperbolic_isGeodesic_of}
  \leanok
  \uses{lem:dc-ch8-3-christoffel-contraction}
  A twice-differentiable curve $u:\mathbb R\to\mathbb R^n$ staying in the
  half-space ($u_n(t)>0$) whose acceleration satisfies the ordinary differential
  equation
  $$
  u''(t)+\frac{1}{u_n(t)}\Big(\langle u'(t),u'(t)\rangle\, e_n-2\,u'_n(t)\,u'(t)\Big)=0
  $$
  is a geodesic of $H^n$. This is the geodesic equation $u''+\Gamma(u',u')=0$ read
  through \cref{lem:dc-ch8-3-christoffel-contraction}; because the chart of the open
  half-space is the inclusion, the moving-foot chart curve of the intrinsic geodesic
  predicate is literally $u$.
\end{lemma}
\begin{proof}
  \leanok
  For an open subset of $\mathbb R^n$ the extended chart at any base point is the
  inclusion, so the chart-local curve $s\mapsto\varphi_{\gamma t}(\gamma s)$ equals
  $u$ for every $t$; its first and second derivatives are $u'$ and $u''$. The
  intrinsic moving-foot geodesic condition at time $t$ is
  $u''(t)+\Gamma(u'(t),u'(t))(u(t))=0$, and substituting the closed form
  \cref{lem:dc-ch8-3-christoffel-contraction} (with $v=w=u'(t)$, so the two mixed
  terms coincide) gives exactly the displayed equation.
\end{proof}

\begin{lemma}[vertical lines are geodesics of $H^n$]
  \label{lem:dc-ch8-3-vertical-geodesic}
  \lean{Riemannian.Hyperbolic.hyperbolic_vertical_isGeodesic}
  \leanok
  \uses{lem:dc-ch8-3-geodesic-equation}
  For a horizontal base point $c$ (with $c_n=0$) and a constant $a>0$, the affinely
  parametrised vertical line
  $$
  t\longmapsto c+(a\,e^{t})\,e_n,\qquad t\in\mathbb R,
  $$
  perpendicular to the boundary hyperplane $\{x_n=0\}$, is a geodesic of $H^n$. Its
  distinguished coordinate $u_n(t)=a\,e^{t}$ solves do Carmo's $u_n''\,u_n=(u_n')^2$.
\end{lemma}
\begin{proof}
  \leanok
  Here $u(t)=c+(a e^{t})e_n$, $u'(t)=u''(t)=(a e^{t})e_n$, $u_n(t)=u'_n(t)=a e^{t}$
  and $\langle u',u'\rangle=(a e^{t})^2$. The geodesic equation
  \cref{lem:dc-ch8-3-geodesic-equation} reads
  $(a e^{t})e_n+(a e^{t})^{-1}\big((a e^{t})^2 e_n-2(a e^{t})^2 e_n\big)
  =(a e^{t})e_n-(a e^{t})e_n=0$, which holds identically.
\end{proof}

\begin{lemma}[semicircles perpendicular to $\partial H^n$ are geodesics of $H^n$]
  \label{lem:dc-ch8-3-semicircle-geodesic}
  \lean{Riemannian.Hyperbolic.hyperbolic_semicircle_isGeodesic}
  \leanok
  \uses{lem:dc-ch8-3-geodesic-equation}
  For a center $m$ on the boundary hyperplane ($m_n=0$), a unit horizontal vector
  $\hat u$ ($\hat u_n=0$, $\langle\hat u,\hat u\rangle=1$) and a radius $r>0$, the
  affinely parametrised semicircle
  $$
  t\longmapsto m+(r\tanh t)\,\hat u+(r\,\mathrm{sech}\,t)\,e_n,\qquad t\in\mathbb R,
  $$
  (with $\tanh t=\sinh t/\cosh t$, $\mathrm{sech}\,t=1/\cosh t$), perpendicular to
  the boundary $\{x_n=0\}$ and centered on it, is a geodesic of $H^n$.
\end{lemma}
\begin{proof}
  \leanok
  With $u(t)=m+(r\tanh t)\hat u+(r\,\mathrm{sech}\,t)e_n$ one has
  $u'(t)=(r\,\mathrm{sech}^2 t)\hat u-(r\tanh t\,\mathrm{sech}\,t)e_n$; using
  $\mathrm{sech}'=-\tanh\,\mathrm{sech}$, $\tanh'=\mathrm{sech}^2$, and
  $\cosh^2-\sinh^2=1$ one computes $u_n=r\,\mathrm{sech}\,t$,
  $\langle u',u'\rangle=r^2\mathrm{sech}^2 t$, and both the $\hat u$- and
  $e_n$-components of $u''+u_n^{-1}(\langle u',u'\rangle e_n-2u'_n u')$ vanish
  identically. Hence \cref{lem:dc-ch8-3-geodesic-equation} applies.
\end{proof}

\begin{proposition}[geodesics of $H^n$]
  \label{prop:dc-ch8-3-1}
  \dcref{ch8:3.1}
  \uses{ex:dc-ch3-3-10,lem:dc-ch8-3-vertical-geodesic,lem:dc-ch8-3-semicircle-geodesic}
  The straight lines perpendicular to the hyperplane $x_n=0$, and the circles of
  $H^n$ whose planes are perpendicular to the hyperplane $x_n=0$ and whose centers
  are in this hyperplane, are the geodesics of $H^n$.
\end{proposition}
\begin{proof}
  An isometry of $\mathbb{R}^n$ which involves solely the variables
  $x_1,\dots,x_{n-1}$ does not change the metric $g_{ij}$, and is therefore an
  isometry of $H^n$. Therefore it suffices to consider lines and circles in the
  $x_1 x_n$-plane. The theorem follows now from \cref{ex:dc-ch3-3-10}. $\square$

  A direct route, independent of \cref{ex:dc-ch3-3-10}, verifies each family
  against the geodesic equation \cref{lem:dc-ch8-3-geodesic-equation}: the vertical
  lines are discharged by \cref{lem:dc-ch8-3-vertical-geodesic} and the semicircles
  perpendicular to $\{x_n=0\}$ by \cref{lem:dc-ch8-3-semicircle-geodesic}. The
  converse inclusion — that these are \emph{all} the geodesics — follows from the
  existence/uniqueness theorem for geodesics, since through every point and
  direction there passes one curve of the stated type.
  
  It follows from the existence and uniqueness theorem for geodesics that all the
  geodesics of $H^n$ are of the type described in Proposition 3.1, and are contained
  in planes perpendicular to the hyperplane $x_n=0$. Since such planes are clearly
  isometric to the hyperbolic plane, the completeness of $H^n$ is a consequence of
  the completeness of the hyperbolic plane (cf. Exercise 10, Chap. 7). Another model
  of hyperbolic space is given in Exercise 3 of this chapter.
\end{proof}

\section{Space forms}

Now we are able to prove the main theorem of this chapter. As always, $M^n$
denotes a manifold of dimension $n$.

\begin{theorem}[universal covering of a space of constant curvature]
  \label{thm:dc-ch8-4-1}
  \dcref{ch8:4.1}
  \notready
  Let $M^n$ be a complete Riemannian manifold with constant sectional curvature $K$.
  Then the universal covering $\tilde{M}$ of $M$, with the covering metric, is
  isometric to:
  - (a) $H^n$, if $K=-1$,
  - (b) $\mathbb{R}^n$, if $K=0$,
  - (c) $S^n$, if $K=1$.
\end{theorem}

\begin{lemma}[uniqueness of a local isometry from a point datum]
  \label{lem:dc-ch8-4-2}
  \dcref{ch8:4.2}
  \uses{thm:dc-ch8-4-1}
  \notready
  Let $f_i:M\to N$, $i=1,2$, be two local isometries of the (connected) Riemannian
  manifold $M$ to the Riemannian manifold $N$. Suppose that there exists a point
  $p\in M$ such that $f_1(p)=f_2(p)$ and $(df_1)_p=(df_2)_p$. Then $f_1=f_2$.
  
  \emph{Proof of the lemma.} Let $V$ be a normal neighborhood of $p$ such that the
  restrictions $f_1|V$ and $f_2|V$ are diffeomorphisms. Let
  $\varphi=f_1^{-1}\circ f_2:V\to V$. Then $\varphi(p)=p$ and $d\varphi_p$ is the
  identity. If $q\in V$, there exists a unique $v\in T_pM$ with $\exp_p(v)=q$, and it
  follows that $\varphi(q)=q$, hence $f_1=f_2$ on $V$. Since $M$ is connected, any
  point $r\in M$ can be joined to $p$ by a path $\alpha:[0,1]\to M$, $\alpha(0)=p$,
  $\alpha(1)=r$. Let
  
  $$
  A=\{t\in[0,1];\ f_1(\alpha(t))=f_2(\alpha(t))\text{ and }(df_1)_{\alpha(t)}=(df_2)_{\alpha(t)}\}.
  $$
  
  From the preceding, $\sup A$ is positive. If $\sup A=t_0\neq 1$, we repeat the
  argument for the point $\alpha(t_0)$, obtaining a contradiction. Therefore
  $\sup A=1$, hence $f_1(r)=f_2(r)$ for all $r\in M$. $\square$
  
  \emph{Proof of \cref{thm:dc-ch8-4-1}.} $\tilde{M}$ is a simply connected, complete Riemannian
  manifold with constant sectional curvature $K$. Consider first the cases (a) and
  (b) and denote, for convenience, both $H^n$ and $\mathbb{R}^n$ by $\triangle$. Fix
  points $p\in\triangle$, $\tilde{p}\in\tilde{M}$ and a linear isometry
  $i:T_p(\triangle)\to T_{\tilde{p}}(\tilde{M})$. Consider the map
  $f=\exp_{\tilde{p}}\circ\,i\circ\exp_p^{-1}:\triangle\to\tilde{M}$. Since $\triangle$
  and $\tilde{M}$ are complete with non-positive sectional curvature, $f$ is
  well-defined. From the Theorem of Cartan, $f$ is a local isometry. From Lemma 3.3
  of Chap. 7, $f$ is a diffeomorphism, which proves (a) and (b).
  
  For case (c), fix $p\in S^n$, $\tilde{p}\in\tilde{M}$ and a linear isometry
  $i:T_p(S^n)\to T_{\tilde{p}}(\tilde{M})$. Let $q\in S^n$ be the antipodal point of
  $p$ and define $f=\exp_{\tilde{p}}\circ\,i\circ\exp_p^{-1}:S^n-\{q\}\to\tilde{M}$.
  From the Theorem of Cartan, $f$ is a local isometry. Choose $p'\in S^n$,
  $p'\neq p$, $p'\neq q$, set $\tilde{p}'=f(p')$, $i'=df_{p'}$ and define
  $f'=\exp_{\tilde{p}'}\circ\,i'\circ\exp_{p'}^{-1}:S^n-\{q'\}\to\tilde{M}$, where
  $q'$ is the antipodal point of $p'$. On the connected set
  $W=S^n-(\{q\}\cup\{q'\})$ we have $p'\in W$ and $f(p')=\tilde{p}'=f'(p')$,
  $df_{p'}=i'=df'_{p'}$, so by Lemma 4.2, $f=f'$ on $W$. We can then define
  $g:S^n\to\tilde{M}$ by $g(r)=f(r)$ if $r\in S^n-\{q\}$ and $g(r)=f'(r)$ if
  $r\in S^n-\{q'\}$. Then $g$ is a local isometry, hence a local diffeomorphism. By
  the compactness of $S^n$, $g$ is a covering map, and since $\tilde{M}$ is simply
  connected, $g$ is a diffeomorphism. Therefore $g$ is an isometry. $\square$
  
  The complete manifolds with constant sectional curvature are called \emph{space forms}.
  The last theorem reduces the determination of all the space forms to a problem in
  group theory. Recall: a group $G$ \emph{acts} on a set $M$ if there is a map
  $G\times M\to M$, $(g,x)\mapsto gx$, with $ex=x$ and $(g_1 g_2)x=g_1(g_2 x)$; $G$
  acts \emph{freely} if $gx=x$ implies $g=e$; the \emph{orbit} of $x$ is $Gx=\{gx;\ g\in G\}$.
  The action is \emph{transitive} if $Gx=M$. If $M$ is a topological space, $G$ (a group
  of homeomorphisms) acts in a \emph{totally discontinuous manner} if every $x\in M$ has a
  neighborhood $U$ such that $g(U)\cap U=\emptyset$ for all $g\in G$, $g\neq e$; in
  this case $\pi:M\to M/G$ is a regular covering map and $G$ is the group of covering
  transformations.
  
  If $M$ is a Riemannian manifold and $\Gamma$ is a subgroup of the group of
  isometries acting in a totally discontinuous manner, then $M/\Gamma$ has a
  differentiable structure in which $\pi:M\to M/\Gamma$ is a local diffeomorphism,
  and a Riemannian metric making $\pi$ a local isometry. Given $p\in M/\Gamma$,
  choose $\tilde p\in\pi^{-1}(p)$; for $u,v\in T_p(M/\Gamma)$, define
  
  $$
  \langle u,v\rangle=\langle d\pi^{-1}(u),d\pi^{-1}(v)\rangle_{\tilde{p}},
  $$
  
  Since $\pi$ is a regular covering, $\Gamma$ is transitive on $\pi^{-1}(p)$, and
  because $\Gamma$ consists of isometries the definition is independent of the
  choice of $\tilde p$. This is called the *metric on $M/\Gamma$ induced by the
  covering $\pi$*. $M/\Gamma$ is complete iff $M$ is, and has constant curvature iff
  $M$ does. Taking $M=S^n$,
  $\mathbb{R}^n$ or $H^n$, we conclude that $M/\Gamma$ is a complete manifold of
  constant curvature $1$, $0$ or $-1$, respectively.
\end{lemma}

\begin{proposition}[every space form is a quotient of a model]
  \label{prop:dc-ch8-4-3}
  \dcref{ch8:4.3}
  Let $M$ be a complete Riemannian manifold with constant sectional curvature $K$
  ($1,0,-1$). Then $M$ is isometric to $\tilde{M}/\Gamma$, where $\tilde{M}$ is $S^n$
  (if $K=1$), $\mathbb{R}^n$ (if $K=0$) or $H^n$ (if $K=-1$), $\Gamma$ is a subgroup
  of the group of isometries of $\tilde{M}$ which acts in a totally discontinuous
  manner on $\tilde{M}$, and the metric on $\tilde{M}/\Gamma$ is induced from the
  covering $\pi:\tilde{M}\to\tilde{M}/\Gamma$.
\end{proposition}
\begin{proof}
  Consider the universal covering $p:\tilde{M}\to M$ with the covering
  metric. Let $\Gamma$ be the group of covering transformations of $p$. Then
  $\Gamma$ is a subgroup of the isometries of $\tilde{M}$ and acts in a totally
  discontinuous manner. Introduce on $\tilde{M}/\Gamma$ the metric induced by
  $\pi:\tilde{M}\to\tilde{M}/\Gamma$. Since the covering $p$ is regular, given
  $\tilde{x},\tilde{y}\in\tilde{M}$, $p(\tilde{x})=p(\tilde{y})$ iff
  $\Gamma\tilde{x}=\Gamma\tilde{y}$, equivalently $\pi(\tilde{x})=\pi(\tilde{y})$.
  The equivalence classes given by $p$ and $\pi$ are the same, which induces a
  bijection $\xi:M\to\tilde{M}/\Gamma$ such that $\pi=\xi\circ p$. Since $\pi$ and
  $p$ are local isometries, $\xi$ is a local isometry, and being a bijection, is an
  isometry of $M$ onto $\tilde{M}/\Gamma$. $\square$
  
  The last proposition reduces the problem of finding all the space forms to that of
  determining all the subgroups of the group of isometries that act in a totally
  discontinuous manner on each of the simply connected models. The spherical problem
  ($\tilde{M}=S^n$) was solved during the sixties (Wolf). Here we mention only two
  interesting facts.
\end{proof}

\begin{proposition}[even-dimensional space forms of $K=1$]
  \label{prop:dc-ch8-4-4}
  \dcref{ch8:4.4}
  \uses{prop:dc-ch8-4-3}
  Let $M^n$ be a complete Riemannian manifold of even dimension $n=2m$, with constant
  sectional curvature $K=1$. Then $M^n$ is isometric to the sphere $S^n$ or the real
  projective space $P^n$ of the same dimension.
\end{proposition}
\begin{proof}
  The orthogonal group $O(2m+1)$ is the (transitive) group of isometries of
  $S^n$. Let $\Gamma$ be a subgroup that acts in a totally discontinuous manner on
  $S^n$. If $\gamma\in\Gamma$ has determinant $+1$, there exists an eigenvalue of
  $\gamma$ equal to $1$; then $\gamma$ has a fixed point, which implies $\gamma=e$.
  If $\tilde{\gamma}\in\Gamma$ has determinant $-1$, then $\tilde{\gamma}^2$ has
  determinant $1$, therefore $\tilde{\gamma}^2=e$. The eigenvalues of
  $\tilde{\gamma}$ are $1$ or $-1$. If $1$ is an eigenvalue then $\tilde{\gamma}=e$,
  contradicting $\det\tilde{\gamma}=-1$. As a result $\tilde{\gamma}=-e$, hence
  $\Gamma=\{e,-e\}$. Using \cref{prop:dc-ch8-4-3} we obtain that $M^n$ is either $S^n$, if
  $\Gamma=\{e\}$, or $P^n$, if $\Gamma=\{e,-e\}$. $\square$
\end{proof}

\begin{proposition}[metric of constant negative curvature on surfaces of genus $>1$]
  \label{prop:dc-ch8-4-5}
  \dcref{ch8:4.5}
  \notready
  Every compact orientable surface of genus $p>1$ can be provided with a metric of
  constant negative curvature.
\end{proposition}
\begin{proof}
  In the hyperbolic plane $H^2$ take a closed geodesic polygon $P$ with $4p$
  sides of equal lengths. By a well known process of identifying the sides, $P$ can
  be identified with a topological surface $M^2$ (Massey, Chap. 1). Let $\Gamma$ be
  the subgroup of isometries of $H^2$ generated by the isometries that identify the
  sides of $P$. The transforms of $P$ by $\Gamma$ "tile $H^2$ without any gaps" iff
  the sum $\alpha$ of the interior angles at the vertices of $P$ equals $2\pi$; in
  this case $M^2=H^2/\Gamma$ has a metric of constant curvature $-1$. We assert that
  it is possible to find a polygon $P$ satisfying $\alpha=2\pi$. By the Gauss–Bonnet
  Theorem, $\alpha=-A+2\pi(2p-1)$, where $A$ is the area of $P$ in the hyperbolic
  metric. If $P$ is arbitrarily small, $\alpha$ is arbitrarily close to $2\pi(2p-1)$;
  if we enlarge the area $A$, $\alpha$ becomes arbitrarily small. By continuously
  deforming $P$, there exists a geodesic polygon such that $\alpha=2\pi$. $\square$
  
  The spaces of constant curvature have an important role in the historical
  development of Riemannian Geometry, due to their relationship with non-Euclidean
  geometry. A non-Euclidean geometry is a complete Riemannian manifold $M$ together
  with a transitive group of isometries $G$ (the non-Euclidean motions) satisfying
  the \emph{Axiom of free mobility}: if $p,\tilde p\in M$, if $\gamma_1,\gamma_2$ are
  geodesics starting at $p$ and forming an angle $\alpha$, and if
  $\tilde\gamma_1,\tilde\gamma_2$ are geodesics starting at $\tilde p$ and forming
  the same angle $\alpha$, then there is a $g\in G$ with $g(p)=\tilde p$,
  $g(\gamma_1)=\tilde\gamma_1$, and $g(\gamma_2)=\tilde\gamma_2$. This corresponds
  to the side-angle-side condition for congruence of triangles in Euclidean geometry
  and implies that $M$ has constant sectional curvature, so the spaces of
  non-Euclidean Geometry are included among the space forms. The cases of $S^n$,
  $P^n$ and $H^n$ are called \emph{spherical}, \emph{elliptic} and \emph{hyperbolic} geometry,
  respectively.
\end{proof}

\section{Isometries of the hyperbolic space; Theorem of Liouville}

The isometries of hyperbolic space are closely related to the conformal
transformations of $\mathbb{R}^n$. A map $f:U\subset\mathbb{R}^n\to\mathbb{R}^n$ of
an open set is called \emph{conformal} if the (non-oriented) angles of intersecting
curves are preserved. The principal objective of this section is to show that the
isometries of the upper half-space $H^n\subset\mathbb{R}^n$ with the metric
$\delta_{ij}/x_n^2$ are the restrictions to $H^n$ of conformal transformations of
$\mathbb{R}^n$ that map $H^n$ onto $H^n$. For $\mathbb{R}^2\cong\mathbb{C}$, the
conformal transformations are the holomorphic or anti-holomorphic functions with
non-zero derivatives. For $\mathbb{R}^n$, $n>2$, the situation is radically
different, and being conformal imposes strong restrictions (Theorem of Liouville).

A necessary and sufficient condition for $f:U\subset\mathbb{R}^n\to\mathbb{R}^n$ to
be conformal is that for all $p\in U$ and all pairs of vectors $v_1,v_2$ at $p$:

$$
\langle df_p(v_1),df_p(v_2)\rangle=\lambda^2(p)\langle v_1,v_2\rangle,\qquad\lambda^2\neq 0.\qquad(3)
$$

Indeed, if $(3)$ holds then $|df_p(v)|^2=\lambda^2|v|^2$, hence

$$
\cos\angle(df_p(v_1),df_p(v_2))=\cos\angle(v_1,v_2),\qquad(4)
$$

so angles are preserved; conversely $(4)$ implies $(3)$. The positive function
$\lambda:U\to\mathbb{R}$ defined in $(3)$ is called the *coefficient of
conformality* of $f$.

\begin{definition}[conformal map with a coefficient of conformality]
  \label{def:dc-ch8-5-conformal}
  \lean{Riemannian.IsConformalWithCoeff,
    Riemannian.IsConformalWithCoeff.isConformalMap,
    Riemannian.IsConformalWithCoeff.conformalAt}
  \leanok
  A map $f:U\subset\mathbb{R}^n\to\mathbb{R}^n$ is \emph{conformal at $p$ with
  coefficient of conformality $\lambda(p)>0$} if $f$ is differentiable at $p$ and
  its differential rescales every inner product by $\lambda^2$, i.e. condition
  $(3)$ holds:
  $\langle df_p(v_1),df_p(v_2)\rangle=\lambda^2(p)\langle v_1,v_2\rangle$ for all
  $v_1,v_2$. Equivalently $df_p$ is a conformal linear map, so $f$ is conformal at
  $p$ in the sense of preserving non-oriented angles.
\end{definition}

\begin{lemma}[a conformal map preserves norms up to $\lambda$ and preserves angles: $(3)\Rightarrow(4)$]
  \label{lem:dc-ch8-5-conformal-angle}
  \uses{def:dc-ch8-5-conformal}
  \lean{Riemannian.IsConformalWithCoeff.norm_fderiv_eq,
    Riemannian.IsConformalWithCoeff.angle_fderiv_eq}
  \leanok
  If $f$ is conformal at $p$ with coefficient $\lambda$, then $|df_p(v)|=\lambda|v|$
  for every $v$, and $df_p$ preserves the non-oriented angle between any two
  vectors: $\angle(df_p(v_1),df_p(v_2))=\angle(v_1,v_2)$, which is do Carmo's
  equation $(4)$.
\end{lemma}
\begin{proof}
  \leanok
  From $(3)$ with $v_1=v_2=v$, $|df_p(v)|^2=\langle df_p(v),df_p(v)\rangle
  =\lambda^2\langle v,v\rangle=\lambda^2|v|^2$, so $|df_p(v)|=\lambda|v|$ since
  $\lambda>0$. Hence the cosine of the angle
  $\langle df_p(v_1),df_p(v_2)\rangle/(|df_p(v_1)|\,|df_p(v_2)|)
  =\lambda^2\langle v_1,v_2\rangle/(\lambda^2|v_1|\,|v_2|)
  =\langle v_1,v_2\rangle/(|v_1|\,|v_2|)$
  is unchanged, giving $(4)$.
\end{proof}

\begin{example}[isometries, dilatations, inversions]
  \label{ex:dc-ch8-5-1}
  \dcref{ch8:5.1}
  \uses{def:dc-ch8-5-conformal}
  \lean{Riemannian.isConformalWithCoeff_isometry,
    Riemannian.isConformalWithCoeff_dilatation,
    Riemannian.isConformalWithCoeff_inversion}
  \leanok
  An isometry of $\mathbb{R}^n$ (orthogonal linear transformation followed by a
  translation) is a conformal transformation with coefficient of conformality
  $\lambda\equiv 1$. The linear transformation $f(p)=\lambda I(p)$, where $I$ is the
  identity matrix and $\lambda=\text{const.}>0$, is conformal with coefficient
  $\lambda$; $f$ is called a \emph{dilatation}. The \emph{inversion} with respect to the unit
  sphere centered at $p_0\in\mathbb{R}^n$, defined by
  
  $$
  f(p)=\frac{p-p_0}{|p-p_0|^2}+p_0,\qquad p\in\mathbb{R}^n-\{p_0\},\qquad(5)
  $$
  
  is a conformal transformation. Geometrically, $f$ takes $p$ to a point on the line
  through $p$ and $p_0$ at distance $1/|p-p_0|$ from $p_0$; $f$ keeps fixed the unit
  sphere around $p_0$ and permutes its interior with its exterior; in particular
  $f^2=\text{identity}$. For a vector $v$ at $p$,
  
  $$
  df_p(v)=\frac{v|p-p_0|^2-2\langle v,p-p_0\rangle(p-p_0)}{|p-p_0|^4},
  $$
  
  and one computes $|df_p(v)|^2=\langle v,v\rangle/|p-p_0|^4$, so the inversion $(5)$
  is conformal with coefficient of conformality $\lambda=1/|p-p_0|^2$.
  
  The Theorem of Liouville asserts the surprising fact that every conformal
  transformation $f$ of an open set $U\subset\mathbb{R}^n$, $n\geq 3$, extends to a
  composition of isometries, dilatations and inversions, at most one of each.
\end{example}

\begin{lemma}[conformal maps are closed under composition]
  \label{lem:dc-ch8-5-conformal-comp}
  \uses{def:dc-ch8-5-conformal}
  \lean{Riemannian.IsConformalWithCoeff.comp,
    Riemannian.isConformalWithCoeff_id}
  \leanok
  If $f$ is conformal at $p$ with coefficient $\lambda$ and $g$ is conformal at
  $f(p)$ with coefficient $\mu$, then $g\circ f$ is conformal at $p$ with
  coefficient $\mu\lambda$. In particular the identity is conformal with
  coefficient $1$, and any composition of isometries, dilatations and inversions —
  the class of maps in the Theorem of Liouville — is conformal.
\end{lemma}
\begin{proof}
  \leanok
  By the chain rule $d(g\circ f)_p=dg_{f(p)}\circ df_p$. For all $v_1,v_2$,
  $\langle d(g\circ f)_p(v_1),d(g\circ f)_p(v_2)\rangle
  =\mu^2\langle df_p(v_1),df_p(v_2)\rangle
  =\mu^2\lambda^2\langle v_1,v_2\rangle=(\mu\lambda)^2\langle v_1,v_2\rangle$,
  and $\mu\lambda>0$.
\end{proof}

\begin{theorem}[Liouville]
  \label{thm:dc-ch8-5-2}
  \dcref{ch8:5.2}
  \notready
  Let $f:U\to\mathbb{R}^n$, $n\geq 3$, be a conformal transformation of an open set
  $U\subset\mathbb{R}^n$. Then $f$ is the restriction to $U$ of a composition of
  isometries, dilatations or inversions, at most one of each.
\end{theorem}
\begin{proof}
  Let $a_1=(1,0,\dots,0),\dots,a_n=(0,\dots,0,1)$ be the canonical basis and
  $(x_1,\dots,x_n)$ the cartesian coordinates. Let $e_1,\dots,e_n$ be parallel
  differentiable vector fields on $U$ with $\langle e_i,e_j\rangle=\delta_{ij}$. If
  $\lambda$ is the coefficient of conformality of $f$,
  
  $$
  \langle df(e_i),df(e_k)\rangle=\lambda^2\delta_{ik},\qquad i,k=1,\dots,n.\qquad(6)
  $$
  
  Let $d^2 f$ be the second differential of $f$, a symmetric bilinear map with
  $d^2 f(a_i,a_j)=\partial^2 f/\partial x_i\partial x_j$. Differentiating $(6)$ with
  $i,j,k$ distinct and combining the three cyclic equations yields
  $\langle d^2 f(e_k,e_j),df(e_i)\rangle=0$ for $i,j,k$ distinct. Letting $i$ vary,
  $d^2 f(e_k,e_j)$ belongs to the plane generated by $df(e_j)$ and $df(e_k)$, so
  $d^2 f(e_k,e_j)=\mu\,df(e_k)+\nu\,df(e_j)$ with $\mu=d\lambda(e_j)/\lambda$,
  $\nu=d\lambda(e_k)/\lambda$, that is,
  
  $$
  d^2 f(e_k,e_j)=\frac{1}{\lambda}(df(e_k)\,d\lambda(e_j)+df(e_j)\,d\lambda(e_k)).\qquad(7)
  $$
  
  Put $\rho=1/\lambda$. Using $d(\rho f)=d\rho\,f+\rho\,df$ and $(7)$,
  
  $$
  d^2(\rho f)(e_k,e_j)=d^2\rho(e_k,e_j)f.\qquad(8)
  $$
  
  We claim $d^2\rho(e_k,e_j)=0$ for $k\neq j$. Computing the third differential
  $d^3(\rho f)$ via $(8)$ and using the symmetry of the first member in $i,j,k$, one
  obtains $d^2\rho(e_k,e_j)df(e_i)=d^2\rho(e_k,e_i)df(e_j)$; since $df(e_i)$ and
  $df(e_j)$ are linearly independent and $i,j,k$ are distinct, $d^2\rho(e_k,e_j)=0$
  for all $j\neq k$. Choosing $e_1,\dots,e_n$ to form a prescribed orthonormal basis
  at any $p$, this holds for every orthonormal basis; since
  $0=d^2\rho\big(\tfrac{e_j+e_k}{\sqrt 2},\tfrac{e_j-e_k}{\sqrt 2}\big)=\tfrac12\{d^2\rho(e_j,e_j)-d^2\rho(e_k,e_k)\}$,
  we get $d^2\rho(e_j,e_j)=d^2\rho(e_k,e_k)$ for $j\neq k$. In summary, for every
  $p\in U$ and any orthonormal basis, $d^2\rho(e_j,e_k)=\sigma\delta_{jk}$; in the
  canonical basis,
  
  $$
  \frac{\partial^2\rho}{\partial x_i\partial x_j}=\sigma\delta_{ij}.\qquad(9)
  $$
  
  Differentiating $(9)$ gives $\partial\sigma/\partial x_i=0$ ($i\neq j$), so
  $\sigma=\text{const.}$
  
  *Case $\sigma=\text{const.}\neq 0$.* Equation $(9)$ implies
  
  $$
  \rho=\frac{\sigma}{2}\sum x_i^2+\sigma\sum b_i x_i+c,\quad b_i,c\text{ constants},\qquad(10)
  $$
  
  which we write as $1/\lambda=\rho=a_1|p-p_0|^2+k_1$ with $a_1=\sigma/2$,
  $k_1=\text{const.}$, $p_0\in\mathbb{R}^n$. The proof in this case will be complete
  once we show $k_1=0$: if $k_1=0$, taking the inversion
  $g(p)=\frac{p-p_0}{|p-p_0|^2}+p_0$ and $h=g\circ f^{-1}$, $h$ is conformal with
  coefficient $a_1$, hence an isometry followed by a dilatation; thus
  $f=h^{-1}\circ g$ is an inversion followed by a dilatation, followed by an
  isometry. To show $k_1=0$, apply the argument to $f^{-1}$:
  $\lambda=a_2|f(p)-q_0|^2+k_2$, hence
  
  $$
  (a_1|p-p_0|^2+k_1)(a_2|f(p)-q_0|^2+k_2)=1.\qquad(11)
  $$
  
  Equation $(11)$ shows that a sphere of center $p_0$ is mapped by $f$ into a sphere
  of center $q_0$, and since $f$ preserves angles, radii map to radii. For a radial
  segment $p(s)$, $0\leq s\leq s_0$, of arc length $s$,
  
  $$
  \int_0^{s_0}\Big|df\Big(\frac{dp}{ds}\Big)\Big|\,ds=\int_0^{s_0}\frac{ds}{a_1|p(s)-p_0|^2+k_1}=|f(p(s_0))-f(p(0))|.
  $$
  
  If $k_1\neq 0$, the left side is a transcendental function of $|p(s_0)-p_0|$, while
  $(11)$ implies it is algebraic. This contradiction shows $k_1=0$, completing the
  case $\sigma\neq0$.
  
  *Case $\sigma=0$.* Here $\rho=1/\lambda=\sum a_i x_i+c_1=A_1(x)+c_1$. Applying the
  initial argument to $f^{-1}$,
  
  $$
  (A_1(x)+c_1)(A_2(f(x))+c_2)=1.\qquad(11')
  $$
  
  Equation $(11')$ shows a hyperplane parallel to $A_1=0$ maps into a hyperplane
  parallel to $A_2=0$; a line perpendicular to $A_1=0$ maps into a line. For a
  segment $p(s)$ of such a line,
  
  $$
  |f(p(s_0))-f(p(0))|=\int_0^{s_0}\frac{ds}{A_1(p(s))+c_1},
  $$
  
  which contradicts $(11')$ unless $A_1(p(s))\equiv 0$. We conclude that if
  $\sigma=0$, $\lambda=\text{const.}$; then lengths of tangent vectors are multiplied
  by a constant $\lambda$ and $f$ is an isometry followed by a dilatation. $\square$
\end{proof}

\begin{theorem}[isometries of $H^n$ are conformal transformations preserving $H^n$]
  \label{thm:dc-ch8-5-3}
  \dcref{ch8:5.3}
  \uses{cor:dc-ch8-2-3}
  The isometries of $H^n$ are the restrictions to $H^n\subset\mathbb{R}^n$ of the
  conformal transformations of $\mathbb{R}^n$ that take $H^n$ onto itself.
\end{theorem}
\begin{proof}
  Suppose first $n\geq 3$. Let $f:H^n\to H^n$ be an isometry in the metric
  $g_{ij}$. By Liouville's Theorem, $f$ extends to a conformal map of $\mathbb{R}^n$
  with the metric $\delta_{ij}$. Since $H^n$ is complete in $g_{ij}$, $f$ maps $H^n$
  onto $H^n$. Conversely, let $f:H^n\to H^n$ be a conformal transformation of $H^n$
  onto $H^n$ and let $e_1,\dots,e_n$ be an orthonormal basis (in $g_{ij}$) at
  $p\in H^n$. Since $g_{ij}=\delta_{ij}/x_n^2$ and $f$ is conformal, there exists
  $\lambda^2>0$ with $\langle df_p(e_i),df_p(e_j)\rangle=\lambda^2\delta_{ij}$, so
  $\{df_p(e_i)/\lambda\}$ is orthonormal at $f(p)$. \cref{cor:dc-ch8-2-3} of Cartan's
  Theorem, together with the global construction after Theorem 2.1 (the exponential
  maps of $H^n$ are global diffeomorphisms), gives an isometry $g$ of $H^n$ taking
  $p$ to $f(p)$ with $dg(e_i)=df(e_i)/\lambda$. By the first part, $g$ is conformal. Hence
  $h=g^{-1}\circ f$ is the restriction to $H^n$ of a conformal map of $\mathbb{R}^n$
  taking $H^n$ onto $H^n$, leaving $p$ fixed and satisfying $dh_p=\lambda I$. It
  remains to show $h=\text{identity}$.
  
  Let $P$ be a hyperplane through $p$. By Liouville's theorem, $h(P)$ is a hyperplane
  or a sphere through $p$; since $dh_p$ is a multiple of the identity, $P$ and $h(P)$
  are tangent at $p$. Since $h$ takes $\partial H^n$ into itself and is conformal,
  the angle of $P$ with $\partial H^n$ equals the angle of $h(P)$ with $\partial H^n$.
  We claim $h(P)=P$. Consider a line $r_1$ through $p$ perpendicular to
  $\partial H^n$, $q_1=r_1\cap\partial H^n$. Since $h(r_1)$ is a circle or a line
  making the same angle with $\partial H^n$, $h(r_1)=r_1$, hence $h(q_1)=q_1$. So if
  $P$ is perpendicular to $\partial H^n$, then $h(P)=P$. If $P$ is not perpendicular,
  let $r$ be a line in $P$; for $h(r)$ to be a circle making the same angle
  $\alpha$ it is necessary that $q_1\in h(r)$ (Fig. 2), contradicting $h(q_1)=q_1$;
  hence $h(r)$ is a line, $h(r)=r$, and $h(P)=P$. It follows that $h$ cannot be an
  inversion (the image of some plane would be a sphere) or an isometry distinct from
  the identity. From Liouville's theorem, $h$ is a dilatation; since $h$ takes $H^n$
  onto $H^n$, $h$ is the identity.
  
  For $n=2$, a simple calculation shows the conformal transformations of the form
  
  $$
  f(z)=\frac{az+b}{cz+d},\quad z\in H^2\subset\mathbb{C},\ a,b,c,d\in\mathbb{R},\ ad-bc>0\qquad(12)
  $$
  
  (which map $H^2$ onto $H^2$) are isometries of $H^2$ with the metric $g_{ij}$.
  Moreover, for a fixed point $p_0\in H^2$ and unit vector $v_0$ at $p_0$, there
  exists a transformation $(12)$ taking an arbitrary $(p,v)$ to $(p_0,v_0)$ ($p$ and
  $v$ are determined by three parameters, the number of parameters of $f$). Since
  there is a unique orientation-preserving isometry of $H^2$ taking $(p,v)$ to
  $(p_0,v_0)$, all orientation-preserving isometries of $H^2$ are of the form (12).
  Composing these with $z\mapsto-\bar z$ gives the orientation-reversing
  anti-holomorphic isometries, so together they give all isometries of $H^2$.
  $\square$
  
  To conclude this section, we identify some important hypersurfaces of $H^n$. The
  intersection with $H^n$ of hyperplanes of $\mathbb{R}^n$ orthogonal to $\partial H^n$,
  and of spheres with center on $\partial H^n$, are totally geodesic submanifolds of
  $H^n$ (Exercise 2). To determine intersections of planes and spheres in any
  position, it is convenient to use the \emph{ball model}: let
  $B^n=\{p\in\mathbb{R}^n;\ |p|<2\}$ with the metric
  
  $$
  h_{ij}(p)=\frac{\delta_{ij}}{(1-\tfrac14|p|^2)^2}.
  $$
\end{proof}

\begin{remark}[the ball model; horospheres and hyperspheres]
  \label{rem:dc-ch8-5-4}
  \dcref{ch8:5.4}
  $B^n$ is isometric to $H^n$ via the map
  
  $$
  f(p)=4\frac{p-p_0}{|p-p_0|^2}-(0,\dots,1),\qquad p_0=(0,\dots,-2),\qquad(13)
  $$
  
  since $\langle df_p(v),df_p(v)\rangle=16\langle v,v\rangle/|p-p_0|^4$ and
  $f_n(p)=(4-|p|^2)/|p-p_0|^2$ give
  $\langle df_p(v),df_p(v)\rangle/(f_n(p))^2=\langle v,v\rangle/(1-\tfrac14|p|^2)^2$.
  Because $f$ is injective, $f$ is an isometry of $B^n$ into $H^n$, taking
  $\partial B^n-\{p_0\}$ into $\partial H^n$. Conjugating by (13) and applying
  \cref{thm:dc-ch8-5-3}, $g:B^n\to B^n$ is an isometry in $h_{ij}$ iff $g$ is the
  restriction of a conformal transformation of $\mathbb{R}^n$ taking $B^n$ onto
  $B^n$.
  
  Let $S\subset H^n$ be an $(n-1)$-sphere of Euclidean space completely contained in
  $H^n$. Then $S$ is a \emph{geodesic sphere} of $H^n$: its image $f^{-1}(S)\subset B^n$ is
  a Euclidean sphere, which can be mapped onto a sphere centered at the origin by an
  isometry of $B^n$; since $h_{ij}$ is symmetric with respect to the origin, the
  result is a geodesic sphere, and so is $S$.
  
  Let $S$ be an $(n-1)$-sphere tangent to $\partial H^n$ at $p$ with
  $S-\{p\}\subset H^n$. By an inversion of $\mathbb{R}^n$ at $p$ (an isometry of
  $H^n$), $S$ is mapped to a hyperplane $P$ parallel to $\partial H^n$. The induced
  metric on $P$ is a multiple of the Euclidean metric, so $P$ has constant curvature
  zero, as does $S-\{p\}$. Such submanifolds are called \emph{horospheres} of $H^n$.
  
  Let $S$ be a Euclidean sphere cutting $\partial H^n$ at an angle $\alpha$, and
  denote $S\cap H^n=\Sigma$. Through an inversion at a point of $S\cap\partial H^n$,
  $\Sigma$ is mapped isometrically into the intersection with $H^n$ of a hyperplane $P$
  cutting $\partial H^n$ at the same angle $\alpha$. Let $Q$ be the hyperplane
  orthogonal to $\partial H^n$ containing $P\cap\partial H^n$. Then $P$ is a
  hypersurface equidistant from the totally geodesic hypersurface $Q$: for a geodesic
  $\gamma_r$ (a semi-circle of radius $r$ with center $0\in P\cap\partial H^n$,
  lying in the plane perpendicular to $P\cap\partial H^n$), a
  homothety with center $0$ takes the circle of radius $r$ into one of any radius, so
  the length of $\gamma_r$ between its intersections with $P$ and $Q$ does not depend
  on $r$. We conclude that $P$, or its isometric image $\Sigma$, is obtained by taking
  geodesics perpendicular to $Q$ and marking a fixed distance. Such hypersurfaces are
  called \emph{equidistant surfaces} (or \emph{hyperspheres}). In Exercise 6 it is shown that
  the hypersurfaces above are exactly the umbilic ones, and their mean and sectional
  curvatures are calculated.
\end{remark}
